\documentclass[aspectratio=169,10pt]{beamer}

% Shared Beamer setup for Fall 2026 course slides.
% Clean Cal Poly-inspired theme: no custom class files or uncommon packages.
\mode<presentation>{
  \usetheme{default}
  \usefonttheme{professionalfonts}
}

\definecolor{CalPolyGreen}{HTML}{154734}
\definecolor{CalPolyGold}{HTML}{BD8B13}
\definecolor{CalPolySage}{HTML}{7FA28F}
\definecolor{CalPolyMint}{HTML}{DDF1E7}
\definecolor{CalPolySand}{HTML}{A29F86}
\definecolor{CalPolyBeige}{HTML}{EFEEDE}
\definecolor{DeckInk}{HTML}{1F2933}
\definecolor{DeckMuted}{HTML}{5F6B7A}
\definecolor{DeckSoft}{HTML}{F7F7F5}

\setbeamersize{text margin left=0.55in,text margin right=0.55in}
\setbeamercolor{normal text}{fg=DeckInk,bg=white}
\setbeamercolor{structure}{fg=CalPolyGreen}
\setbeamercolor{title}{fg=white}
\setbeamercolor{frametitle}{fg=CalPolyGreen,bg=white}
\setbeamercolor{block title}{bg=CalPolySage,fg=white}
\setbeamercolor{block body}{bg=CalPolyMint,fg=black}
\setbeamercolor{alerted text}{fg=CalPolyGold}

\setbeamerfont{title}{size=\Huge,series=\bfseries}
\setbeamerfont{frametitle}{size=\Large,series=\bfseries}
\setbeamerfont{block title}{size=\normalsize,series=\bfseries}

\setbeamertemplate{navigation symbols}{}
\setbeamertemplate{itemize item}{\color{CalPolyGold}$\blacktriangleright$}
\setbeamertemplate{itemize subitem}{\color{CalPolyGreen}$\bullet$}
\setbeamertemplate{footline}{
  \hfill{\color{DeckMuted}\scriptsize\insertframenumber/\inserttotalframenumber}
  \hspace{0.18in}\vspace{0.12in}
}
\setbeamertemplate{frametitle}{
  \vspace{0.18in}
  \hspace{0.02in}\insertframetitle\par
  \vspace{0.08in}
  \noindent\makebox[\linewidth][l]{%
    {\color{CalPolyGold}\rule{0.55in}{2pt}}%
    {\color{CalPolyGreen}\rule{\dimexpr\linewidth-0.55in\relax}{2pt}}%
  }
}

\newcommand{\CourseNumber}{Course}
\newcommand{\CourseTitle}{Course Title}
\newcommand{\LectureNumber}{0}
\newcommand{\LectureTitle}{Lecture Title}
\newcommand{\LectureDate}{Date}
\newcommand{\CourseUnit}{Unit}

\newcommand{\coursenumber}[1]{\renewcommand{\CourseNumber}{#1}}
\newcommand{\coursetitle}[1]{\renewcommand{\CourseTitle}{#1}}
\newcommand{\lecturenumber}[1]{\renewcommand{\LectureNumber}{#1}}
\newcommand{\lecturetitle}[1]{\renewcommand{\LectureTitle}{#1}}
\newcommand{\lecturedate}[1]{\renewcommand{\LectureDate}{#1}}
\newcommand{\courseunit}[1]{\renewcommand{\CourseUnit}{#1}}

\author{Kirk Duran}
\institute{California Polytechnic State University, San Luis Obispo}
\date{\LectureDate}

\newcommand{\makedecktitle}{
  \begingroup
  \setbeamercolor{background canvas}{bg=CalPolyGreen}
  \begin{frame}[plain]
    \vfill
    \begin{center}
      {\usebeamerfont{title}\color{white}\LectureTitle\par}
      \vspace{0.18in}
      {\Large\color{white}Lecture \LectureNumber\par}
    \end{center}
    \vfill
    \noindent{\color{white}\rule{\linewidth}{1.2pt}}\par
    \vspace{0.08in}
    \noindent\colorbox{CalPolyGold}{\makebox[0.24\linewidth][c]{\color{white}\Large\CourseNumber}}
    \hspace{0.12in}{\color{white}\Large Kirk Duran}
  \end{frame}
  \endgroup
}

\newenvironment{keyidea}{\begin{block}{Key idea}}{\end{block}}
\newenvironment{activity}{\begin{block}{In class}}{\end{block}}
\newenvironment{cpdefinition}{\begin{block}{Definition}}{\end{block}}
\newenvironment{exampleblockcp}{
  \begingroup
  \setbeamercolor{block title}{bg=CalPolySand,fg=white}
  \setbeamercolor{block body}{bg=CalPolyBeige,fg=black}
  \begin{block}{Example}
}{
  \end{block}
  \endgroup
}

\newcommand{\imageplaceholder}[2][0.3in]{%
  \noindent\makebox[\linewidth][c]{%
    \fcolorbox{CalPolySage}{CalPolyBeige}{%
      \parbox[c][#1][c]{0.86\linewidth}{%
        \centering
        {\color{DeckMuted}\scriptsize Image placeholder: }%
        {\color{DeckInk}\scriptsize\ttfamily\detokenize{#2}}%
      }%
    }%
  }%
}

\newcommand{\sectionframe}[1]{
  \begingroup
  \setbeamercolor{background canvas}{bg=CalPolyGreen}
  \begin{frame}[plain]
    \vfill
    {\color{white}\LARGE\bfseries #1\par}
    \vspace{0.18in}
    {\color{CalPolyGold}\rule{0.22\paperwidth}{3pt}}
    {\color{white}\rule{0.62\paperwidth}{3pt}}
    \vfill
  \end{frame}
  \endgroup
}

\newcommand{\placeholdernote}{
  \begin{block}{Placeholder}
    Replace these bullets with the final lecture material, examples, and in-class activities.
  \end{block}
}


\pdfstringdefDisableCommands{\def\translate#1{#1}}

\graphicspath{{../assets/lecture-13/}}

% If an image file is not yet available, the slide displays a labeled
% placeholder using the intended filename. Place the matching file in
% ../assets/lecture-13/ to replace the placeholder automatically.
\newcommand{\lectureimage}[2][1.75in]{%
  \IfFileExists{../assets/lecture-13/#2}{%
    \includegraphics[width=\linewidth,height=#1,keepaspectratio]{#2}%
  }{%
    \fbox{%
      \parbox[c][#1][c]{0.94\linewidth}{%
        \centering
        \small\textbf{Image Placeholder}\\[0.08in]
        \scriptsize\texttt{\detokenize{#2}}
      }%
    }%
  }%
}

\setbeamersize{text margin left=0.42in,text margin right=0.42in}

\coursenumber{CS 1001}
\coursetitle{Introduction to Computing}
\lecturenumber{13}
\lecturetitle{Tracing Conditionals and Return}
\lecturedate{September 23, 2026}
\courseunit{1. Computing and Program Execution}

\begin{document}

\makedecktitle

\begin{frame}{Today}
  \begin{itemize}
    \item Trace conditional control flow inside a function call.
    \item Explain why a function may contain several \texttt{return} statements but one call follows only one return path.
    \item Analyze chained and nested conditionals that return different values.
    \item Use early returns to terminate a call as soon as its result is known.
    \item Trace a function that calls a predicate function in its condition.
    \item Identify fall-through paths, implicit \texttt{None}, and unreachable statements.
  \end{itemize}

  \begin{keyidea}
    A conditional chooses a path; \texttt{return} terminates the current function call and sends one value back to its caller.
  \end{keyidea}
\end{frame}

\sectionframe{Part 1: Conditional Paths Inside a Function Call}

\begin{frame}{Two Control-Flow Mechanisms Now Interact}
  \begin{columns}[T,onlytextwidth]
    \begin{column}{0.54\textwidth}
      \begin{itemize}
        \item A conditional decides which suite executes next.
        \item A function call establishes a local execution environment.
        \item A \texttt{return} statement ends that call immediately.
      \end{itemize}

      \begin{block}{Example}
        \texttt{def sign(x):}\\
        \hspace*{0.20in}\texttt{if x < 0:}\\
        \hspace*{0.40in}\texttt{return "negative"}\\
        \hspace*{0.20in}\texttt{return "nonnegative"}
      \end{block}
    \end{column}

    \begin{column}{0.40\textwidth}
      % Intended visual: function frame containing a branch with two exit arrows.
      \lectureimage[2.55in]{conditional-return-control-flow.png}
    \end{column}
  \end{columns}
\end{frame}

\begin{frame}{\texttt{return} Is a Control-Flow Statement}
  \begin{cpdefinition}
    Executing \texttt{return expression} evaluates the expression, terminates the current function call, and makes the resulting value the value of that call.
  \end{cpdefinition}

  \begin{block}{Example}
    \texttt{def absolute\_value(x):}\\
    \hspace*{0.20in}\texttt{if x < 0:}\\
    \hspace*{0.40in}\texttt{return -x}\\
    \hspace*{0.20in}\texttt{return x}
  \end{block}

  \begin{keyidea}
    \texttt{return} does not merely leave the current \texttt{if} suite; it ends the entire active function call.
  \end{keyidea}
\end{frame}

\begin{frame}{Multiple \texttt{return} Statements Are Alternative Exit Points}
  \begin{columns}[T,onlytextwidth]
    \begin{column}{0.47\textwidth}
      \begin{block}{Source code}
        \texttt{def classify(n):}\\
        \hspace*{0.20in}\texttt{if n < 0:}\\
        \hspace*{0.40in}\texttt{return "negative"}\\
        \hspace*{0.20in}\texttt{if n == 0:}\\
        \hspace*{0.40in}\texttt{return "zero"}\\
        \hspace*{0.20in}\texttt{return "positive"}
      \end{block}
    \end{column}

    \begin{column}{0.47\textwidth}
      \begin{itemize}
        \item The definition contains three syntactic \texttt{return} statements.
        \item A particular call reaches at most one of them.
        \item Once one executes, the remaining body is not evaluated during that call.
      \end{itemize}
    \end{column}
  \end{columns}
\end{frame}

\begin{frame}{Trace Example: Conditional Return}
  \begin{block}{Program}
    \texttt{def sign(x):}\\
    \hspace*{0.20in}\texttt{if x < 0:}\\
    \hspace*{0.40in}\texttt{return -1}\\
    \hspace*{0.20in}\texttt{return 1}\\[0.08in]
    \texttt{result = sign(-4)}
  \end{block}

  \begin{center}
    \renewcommand{\arraystretch}{1.18}
    \begin{tabular}{c|l|l}
      \textbf{Step} & \textbf{Execution} & \textbf{State / result} \\\hline
      1 & call \texttt{sign(-4)} & \texttt{x -> -4} \\
      2 & evaluate \texttt{x < 0} & \texttt{True} \\
      3 & execute \texttt{return -1} & call produces \texttt{-1} \\
      4 & resume caller & \texttt{result -> -1}
    \end{tabular}
  \end{center}
\end{frame}

\begin{frame}{Statements After an Executed \texttt{return} Are Skipped}
  \begin{block}{Program}
    \texttt{def example(x):}\\
    \hspace*{0.20in}\texttt{if x > 0:}\\
    \hspace*{0.40in}\texttt{return x * 2}\\
    \hspace*{0.40in}\texttt{print("unreachable on this path")}\\
    \hspace*{0.20in}\texttt{return 0}
  \end{block}

  \begin{itemize}
    \item For \texttt{x > 0}, the first \texttt{return} ends the call.
    \item The following \texttt{print} cannot execute on that path.
    \item For \texttt{x <= 0}, execution skips the suite and reaches \texttt{return 0}.
  \end{itemize}

  \begin{alertblock}{Tracing rule}
    After an executed \texttt{return}, immediately remove the active call from the trace and resume the caller.
  \end{alertblock}
\end{frame}

\sectionframe{Part 2: Chained Conditionals with Return Values}

\begin{frame}{A Chained Conditional Can Select a Return Value}
  \begin{block}{Example}
    \texttt{def letter\_grade(score):}\\
    \hspace*{0.20in}\texttt{if score >= 90:}\\
    \hspace*{0.40in}\texttt{return "A"}\\
    \hspace*{0.20in}\texttt{elif score >= 80:}\\
    \hspace*{0.40in}\texttt{return "B"}\\
    \hspace*{0.20in}\texttt{elif score >= 70:}\\
    \hspace*{0.40in}\texttt{return "C"}\\
    \hspace*{0.20in}\texttt{else:}\\
    \hspace*{0.40in}\texttt{return "below C"}
  \end{block}

  \begin{keyidea}
    The chain chooses one suite; that suite's \texttt{return} determines the value of the call.
  \end{keyidea}
\end{frame}

\begin{frame}{Trace a Chained Return Top to Bottom}
  \begin{block}{Call}
    \centering
    \Large\texttt{letter\_grade(84)}
  \end{block}

  \begin{enumerate}
    \item Bind \texttt{score -> 84}.
    \item Evaluate \texttt{score >= 90}: \texttt{False}.
    \item Evaluate \texttt{score >= 80}: \texttt{True}.
    \item Execute \texttt{return "B"}.
    \item Terminate the call; the later tests are not evaluated.
  \end{enumerate}

  \begin{block}{Result}
    \centering
    \texttt{letter\_grade(84)} $\longrightarrow$ \texttt{"B"}
  \end{block}
\end{frame}

\begin{frame}{Order Still Matters in a Chained Return}
  \begin{columns}[T,onlytextwidth]
    \begin{column}{0.47\textwidth}
      \begin{block}{Correct ordering}
        \texttt{if score >= 90:}\\
        \hspace*{0.20in}\texttt{return "A"}\\
        \texttt{elif score >= 80:}\\
        \hspace*{0.20in}\texttt{return "B"}
      \end{block}
    \end{column}

    \begin{column}{0.47\textwidth}
      \begin{block}{Incorrect ordering}
        \texttt{if score >= 80:}\\
        \hspace*{0.20in}\texttt{return "B"}\\
        \texttt{elif score >= 90:}\\
        \hspace*{0.20in}\texttt{return "A"}
      \end{block}
    \end{column}
  \end{columns}

  \begin{itemize}
    \item A score of \texttt{95} satisfies both predicates.
    \item An \texttt{elif} chain selects the first truthy branch.
    \item Later branches cannot repair a classification already returned.
  \end{itemize}
\end{frame}

\begin{frame}{An Exhaustive Chain Guarantees an Explicit Return}
  \begin{columns}[T,onlytextwidth]
    \begin{column}{0.48\textwidth}
      \begin{block}{Exhaustive}
        \texttt{if n < 0:}\\
        \hspace*{0.20in}\texttt{return "negative"}\\
        \texttt{elif n == 0:}\\
        \hspace*{0.20in}\texttt{return "zero"}\\
        \texttt{else:}\\
        \hspace*{0.20in}\texttt{return "positive"}
      \end{block}
    \end{column}

    \begin{column}{0.46\textwidth}
      \begin{itemize}
        \item Every possible numeric input selects one branch.
        \item Every selected branch contains a \texttt{return}.
        \item Therefore every call reaches an explicit return value.
      \end{itemize}
    \end{column}
  \end{columns}
\end{frame}

\begin{frame}{Independent \texttt{if}s with Early Returns Behave Differently}
  \begin{block}{Example}
    \texttt{def classify(n):}\\
    \hspace*{0.20in}\texttt{if n < 0:}\\
    \hspace*{0.40in}\texttt{return "negative"}\\
    \hspace*{0.20in}\texttt{if n == 0:}\\
    \hspace*{0.40in}\texttt{return "zero"}\\
    \hspace*{0.20in}\texttt{return "positive"}
  \end{block}

  \begin{itemize}
    \item If the first condition is false, execution continues to the second \texttt{if}.
    \item If either early return executes, the later tests are never reached.
    \item The final return handles the remaining path.
  \end{itemize}

  \begin{keyidea}
    Early returns can make several independent tests behave like mutually exclusive exit paths.
  \end{keyidea}
\end{frame}

\begin{frame}{Prediction: Which Return Executes?}
  \begin{block}{Function}
    \texttt{def shipping(weight):}\\
    \hspace*{0.20in}\texttt{if weight <= 0:}\\
    \hspace*{0.40in}\texttt{return "invalid"}\\
    \hspace*{0.20in}\texttt{elif weight <= 5:}\\
    \hspace*{0.40in}\texttt{return "standard"}\\
    \hspace*{0.20in}\texttt{elif weight <= 20:}\\
    \hspace*{0.40in}\texttt{return "heavy"}\\
    \hspace*{0.20in}\texttt{return "oversize"}
  \end{block}

  \begin{activity}
    Predict the returned value for \texttt{-2}, \texttt{5}, \texttt{12}, and \texttt{25}. For each call, identify exactly which conditions are evaluated.
  \end{activity}
\end{frame}

\sectionframe{Part 3: Nested Conditionals and Early Returns}

\begin{frame}{A Return Inside a Nested Conditional Ends the Whole Call}
  \begin{block}{Example}
    \texttt{def admission(age, has\_id):}\\
    \hspace*{0.20in}\texttt{if age >= 18:}\\
    \hspace*{0.40in}\texttt{if has\_id:}\\
    \hspace*{0.60in}\texttt{return "admit"}\\
    \hspace*{0.40in}\texttt{return "need id"}\\
    \hspace*{0.20in}\texttt{return "underage"}
  \end{block}

  \begin{itemize}
    \item The outer condition is evaluated first.
    \item The inner condition is evaluated only when the outer condition is true.
    \item Any executed \texttt{return} terminates \texttt{admission}, not merely the nearest conditional.
  \end{itemize}
\end{frame}

\begin{frame}{Trace Nested Returns from the Outside In}
  \begin{block}{Call}
    \centering
    \Large\texttt{admission(21, False)}
  \end{block}

  \begin{center}
    \renewcommand{\arraystretch}{1.18}
    \begin{tabular}{c|l|l}
      \textbf{Step} & \textbf{Evaluation} & \textbf{Outcome} \\\hline
      1 & bind parameters & \texttt{age -> 21, has\_id -> False} \\
      2 & \texttt{age >= 18} & \texttt{True} \\
      3 & \texttt{has\_id} & \texttt{False} \\
      4 & \texttt{return "need id"} & call terminates
    \end{tabular}
  \end{center}

  \begin{block}{Returned value}
    \centering
    \texttt{"need id"}
  \end{block}
\end{frame}

\begin{frame}{Nested Structure Is Appropriate for Dependent Questions}
  \begin{columns}[T,onlytextwidth]
    \begin{column}{0.53\textwidth}
      \begin{itemize}
        \item Sometimes a second question is meaningful only after the first condition succeeds.
        \item Nesting represents this dependency directly.
        \item The inner condition is unreachable when the outer condition is false.
      \end{itemize}

      \begin{block}{Logical structure}
        Is the account active?\\
        \hspace*{0.20in}$\downarrow$ yes\\
        Has the account completed verification?
      \end{block}
    \end{column}

    \begin{column}{0.41\textwidth}
      % Intended visual: hierarchical decision tree with inner question only on true branch.
      \lectureimage[2.45in]{nested-return-decision-tree.png}
    \end{column}
  \end{columns}
\end{frame}

\begin{frame}{Early Returns Can Reduce Unnecessary Nesting}
  \begin{columns}[T,onlytextwidth]
    \begin{column}{0.47\textwidth}
      \begin{block}{Nested form}
        \texttt{if amount >= 0:}\\
        \hspace*{0.20in}\texttt{if amount <= limit:}\\
        \hspace*{0.40in}\texttt{return amount}\\
        \hspace*{0.20in}\texttt{else:}\\
        \hspace*{0.40in}\texttt{return limit}\\
        \texttt{else:}\\
        \hspace*{0.20in}\texttt{return 0}
      \end{block}
    \end{column}

    \begin{column}{0.47\textwidth}
      \begin{block}{Early-return form}
        \texttt{if amount < 0:}\\
        \hspace*{0.20in}\texttt{return 0}\\
        \texttt{if amount > limit:}\\
        \hspace*{0.20in}\texttt{return limit}\\
        \texttt{return amount}
      \end{block}
    \end{column}
  \end{columns}

  \begin{keyidea}
    Equivalent behavior can sometimes be expressed with shallower control flow once invalid or exceptional cases return early.
  \end{keyidea}
\end{frame}

\begin{frame}{A Guard Clause Rejects a Case Before the Main Path}
  \begin{cpdefinition}
    A \textbf{guard clause} is an early conditional return that handles a case before the main computation continues.
  \end{cpdefinition}

  \begin{block}{Example}
    \texttt{def ratio(total, count):}\\
    \hspace*{0.20in}\texttt{if count == 0:}\\
    \hspace*{0.40in}\texttt{return None}\\
    \hspace*{0.20in}\texttt{result = total / count}\\
    \hspace*{0.20in}\texttt{return result}
  \end{block}

  \begin{itemize}
    \item The guard handles the invalid denominator.
    \item The division is evaluated only on the remaining path.
  \end{itemize}
\end{frame}

\begin{frame}{Trace Example: Guard Clause}
  \begin{columns}[T,onlytextwidth]
    \begin{column}{0.47\textwidth}
      \begin{block}{Call A}
        \texttt{ratio(10, 0)}\\[0.08in]
        \texttt{count == 0} $\to$ \texttt{True}\\
        \texttt{return None}\\[0.08in]
        Division is never evaluated.
      \end{block}
    \end{column}

    \begin{column}{0.47\textwidth}
      \begin{block}{Call B}
        \texttt{ratio(10, 2)}\\[0.08in]
        \texttt{count == 0} $\to$ \texttt{False}\\
        \texttt{result -> 5.0}\\
        \texttt{return 5.0}
      \end{block}
    \end{column}
  \end{columns}
\end{frame}

\sectionframe{Part 4: Predicate Calls as Conditions}

\begin{frame}{A Predicate Call Can Determine the Branch}
  \begin{cpdefinition}
    A \textbf{predicate function} returns a Boolean value representing a claim about its arguments.
  \end{cpdefinition}

  \begin{block}{Example}
    \texttt{def is\_even(n):}\\
    \hspace*{0.20in}\texttt{return n \% 2 == 0}\\[0.08in]
    \texttt{def describe(n):}\\
    \hspace*{0.20in}\texttt{if is\_even(n):}\\
    \hspace*{0.40in}\texttt{return "even"}\\
    \hspace*{0.20in}\texttt{return "odd"}
  \end{block}

  \begin{keyidea}
    The predicate call must finish and produce a truth value before the surrounding \texttt{if} can select a branch.
  \end{keyidea}
\end{frame}

\begin{frame}{Trace the Predicate Call Before the Conditional Branch}
  \begin{block}{Call}
    \centering
    \Large\texttt{describe(7)}
  \end{block}

  \begin{enumerate}
    \item Enter \texttt{describe}; bind \texttt{n -> 7}.
    \item Evaluate the condition \texttt{is\_even(n)}.
    \item Call \texttt{is\_even(7)}; its local \texttt{n -> 7}.
    \item Evaluate \texttt{7 \% 2 == 0} $\to$ \texttt{False}.
    \item Return \texttt{False} to \texttt{describe}.
    \item The \texttt{if} condition is false; execute \texttt{return "odd"}.
  \end{enumerate}
\end{frame}

\begin{frame}{The Two Calls Have Separate Parameter Bindings}
  \begin{columns}[T,onlytextwidth]
    \begin{column}{0.45\textwidth}
      \begin{block}{\texttt{describe} frame}
        \texttt{n -> 7}\\[0.08in]
        waiting for condition value while \texttt{is\_even} executes
      \end{block}
    \end{column}

    \begin{column}{0.45\textwidth}
      \begin{block}{\texttt{is\_even} frame}
        \texttt{n -> 7}\\[0.08in]
        evaluates predicate and returns \texttt{False}
      \end{block}
    \end{column}
  \end{columns}

  \begin{keyidea}
    Identical parameter names in separate active calls represent distinct bindings in distinct call frames.
  \end{keyidea}
\end{frame}

\begin{frame}{A Predicate Can Hide a More Detailed Test}
  \begin{block}{Predicate}
    \texttt{def is\_valid\_score(score):}\\
    \hspace*{0.20in}\texttt{return score >= 0 and score <= 100}
  \end{block}

  \begin{block}{Classification function}
    \texttt{def result\_label(score):}\\
    \hspace*{0.20in}\texttt{if not is\_valid\_score(score):}\\
    \hspace*{0.40in}\texttt{return "invalid"}\\
    \hspace*{0.20in}\texttt{if score >= 70:}\\
    \hspace*{0.40in}\texttt{return "passing"}\\
    \hspace*{0.20in}\texttt{return "not passing"}
  \end{block}
\end{frame}

\begin{frame}{Trace Across the Predicate Boundary}
  \begin{block}{Call: \texttt{result\_label(105)}}
    \begin{enumerate}
      \item \texttt{score -> 105} in \texttt{result\_label}.
      \item Evaluate \texttt{is\_valid\_score(score)}.
      \item Predicate returns \texttt{False}.
      \item \texttt{not False} evaluates to \texttt{True}.
      \item Execute \texttt{return "invalid"}.
      \item No later condition in \texttt{result\_label} is evaluated.
    \end{enumerate}
  \end{block}

  \begin{activity}
    Repeat the trace for \texttt{result\_label(82)} and identify both returned values: first from the predicate, then from the classification function.
  \end{activity}
\end{frame}

\begin{frame}{Predicate Results Participate in Ordinary Boolean Expressions}
  \begin{block}{Example}
    \texttt{if is\_valid\_score(score) and score >= 90:}\\
    \hspace*{0.20in}\texttt{return "high valid score"}
  \end{block}

  \begin{itemize}
    \item The predicate call is an expression whose value is Boolean.
    \item The larger condition follows the ordinary rules for Boolean expressions.
    \item To trace accurately, reduce function-call expressions to returned values before deciding the branch.
  \end{itemize}

  \begin{block}{Expression reduction}
    \centering
    \texttt{is\_valid\_score(95) and 95 >= 90}\\
    $\longrightarrow$ \texttt{True and True}\\
    $\longrightarrow$ \texttt{True}
  \end{block}
\end{frame}

\sectionframe{Part 5: Fall-Through, Implicit Return, and Trace Errors}

\begin{frame}{Not Every Conditional Function Returns on Every Path}
  \begin{block}{Problematic definition}
    \texttt{def category(n):}\\
    \hspace*{0.20in}\texttt{if n > 0:}\\
    \hspace*{0.40in}\texttt{return "positive"}
  \end{block}

  \begin{itemize}
    \item \texttt{category(5)} reaches the explicit return.
    \item \texttt{category(0)} skips the suite and reaches the end of the function body.
    \item Reaching the end without executing \texttt{return} produces \texttt{None}.
  \end{itemize}

  \begin{alertblock}{Tracing question}
    For every possible path through a value-producing function, ask: \emph{where does this path return?}
  \end{alertblock}
\end{frame}

\begin{frame}{Trace Fall-Through Explicitly}
  \begin{block}{Call}
    \centering
    \Large\texttt{answer = category(0)}
  \end{block}

  \begin{center}
    \renewcommand{\arraystretch}{1.18}
    \begin{tabular}{c|l|l}
      \textbf{Step} & \textbf{Execution} & \textbf{Result} \\\hline
      1 & enter \texttt{category} & \texttt{n -> 0} \\
      2 & evaluate \texttt{n > 0} & \texttt{False} \\
      3 & reach end of function & implicit \texttt{None} \\
      4 & resume caller & \texttt{answer -> None}
    \end{tabular}
  \end{center}
\end{frame}

\begin{frame}{\texttt{print} Still Does Not Replace \texttt{return}}
  \begin{columns}[T,onlytextwidth]
    \begin{column}{0.47\textwidth}
      \begin{block}{Prints a classification}
        \texttt{def classify(n):}\\
        \hspace*{0.20in}\texttt{if n > 0:}\\
        \hspace*{0.40in}\texttt{print("positive")}\\
        \hspace*{0.20in}\texttt{else:}\\
        \hspace*{0.40in}\texttt{print("nonpositive")}
      \end{block}
    \end{column}

    \begin{column}{0.47\textwidth}
      \begin{block}{Returns a classification}
        \texttt{def classify(n):}\\
        \hspace*{0.20in}\texttt{if n > 0:}\\
        \hspace*{0.40in}\texttt{return "positive"}\\
        \hspace*{0.20in}\texttt{else:}\\
        \hspace*{0.40in}\texttt{return "nonpositive"}
      \end{block}
    \end{column}
  \end{columns}

  \begin{keyidea}
    Conditional output and conditional return values are different program behaviors.
  \end{keyidea}
\end{frame}

\begin{frame}{Unreachable Code Can Be Path-Specific or Unconditional}
  \begin{columns}[T,onlytextwidth]
    \begin{column}{0.47\textwidth}
      \begin{block}{Path-specific}
        \texttt{if x > 0:}\\
        \hspace*{0.20in}\texttt{return x}\\
        \hspace*{0.20in}\texttt{print("never here")}\\
        \texttt{return 0}
      \end{block}

      The \texttt{print} is unreachable only after the true branch is selected.
    \end{column}

    \begin{column}{0.47\textwidth}
      \begin{block}{Unconditionally unreachable}
        \texttt{return x}\\
        \texttt{print("never")}
      \end{block}

      Once the first statement executes, no path reaches the second statement.
    \end{column}
  \end{columns}
\end{frame}

\begin{frame}{Common Tracing Error: Continuing After Return}
  \begin{block}{Program}
    \texttt{def fee(age):}\\
    \hspace*{0.20in}\texttt{if age < 5:}\\
    \hspace*{0.40in}\texttt{return 0}\\
    \hspace*{0.20in}\texttt{if age < 18:}\\
    \hspace*{0.40in}\texttt{return 8}\\
    \hspace*{0.20in}\texttt{return 12}
  \end{block}

  \begin{activity}
    Trace \texttt{fee(3)}. Which conditions are actually evaluated? Then trace \texttt{fee(10)} and \texttt{fee(25)}.
  \end{activity}

  \begin{alertblock}{Rule}
    Never continue tracing the body of a call after one of its \texttt{return} statements has executed.
  \end{alertblock}
\end{frame}

\sectionframe{Part 6: Integrated Tracing}

\begin{frame}{Integrated Trace: Chaining Inside a Function}
  \begin{block}{Program}
    \texttt{def shipping\_cost(weight):}\\
    \hspace*{0.20in}\texttt{if weight <= 0:}\\
    \hspace*{0.40in}\texttt{return None}\\
    \hspace*{0.20in}\texttt{elif weight <= 5:}\\
    \hspace*{0.40in}\texttt{return 6.0}\\
    \hspace*{0.20in}\texttt{elif weight <= 20:}\\
    \hspace*{0.40in}\texttt{return 12.0}\\
    \hspace*{0.20in}\texttt{return 20.0}\\[0.08in]
    \texttt{cost = shipping\_cost(14)}
  \end{block}

  \begin{activity}
    Produce the complete trace: parameter binding, every evaluated condition, selected branch, returned value, and final caller binding.
  \end{activity}
\end{frame}

\begin{frame}{Integrated Trace: Predicate Call in the Condition}
  \begin{block}{Program}
    \texttt{def is\_adult(age):}\\
    \hspace*{0.20in}\texttt{return age >= 18}\\[0.06in]
    \texttt{def ticket\_price(age):}\\
    \hspace*{0.20in}\texttt{if is\_adult(age):}\\
    \hspace*{0.40in}\texttt{return 15}\\
    \hspace*{0.20in}\texttt{return 10}\\[0.06in]
    \texttt{price = ticket\_price(16)}
  \end{block}

  \begin{itemize}
    \item Trace both active call frames when \texttt{is\_adult} executes.
    \item Record the predicate's returned value separately from \texttt{ticket\_price}'s returned value.
    \item Finish by recording the caller binding for \texttt{price}.
  \end{itemize}
\end{frame}

\begin{frame}[shrink=4]{Integrated Trace: Nested Conditionals with Multiple Returns}
  \begin{block}{Program}
    \texttt{def access\_level(active, admin, verified):}\\
    \hspace*{0.20in}\texttt{if active:}\\
    \hspace*{0.40in}\texttt{if admin:}\\
    \hspace*{0.60in}\texttt{return "administrator"}\\
    \hspace*{0.40in}\texttt{elif verified:}\\
    \hspace*{0.60in}\texttt{return "member"}\\
    \hspace*{0.40in}\texttt{else:}\\
    \hspace*{0.60in}\texttt{return "limited"}\\
    \hspace*{0.20in}\texttt{return "inactive"}
  \end{block}

  \begin{activity}
    Trace the conditions and return path for \texttt{access\_level(True, False, True)} and \texttt{access\_level(False, True, True)}. Do not evaluate conditions that execution never reaches.
  \end{activity}
\end{frame}

\begin{frame}{A Reliable Manual Tracing Procedure}
  \begin{enumerate}
    \item Enter the function and record parameter bindings.
    \item Execute statements in order.
    \item When a condition is reached, evaluate its expression completely.
    \item Follow only the selected branch.
    \item If a function call appears inside the condition, trace that call to its returned value first.
    \item When \texttt{return} executes, evaluate its expression and terminate the current call immediately.
    \item Resume the caller using the returned value.
  \end{enumerate}
\end{frame}

\begin{frame}{Final Prediction}
  \begin{block}{Program}
    \texttt{def in\_range(x):}\\
    \hspace*{0.20in}\texttt{return x >= 0 and x <= 10}\\[0.06in]
    \texttt{def transform(x):}\\
    \hspace*{0.20in}\texttt{if not in\_range(x):}\\
    \hspace*{0.40in}\texttt{return None}\\
    \hspace*{0.20in}\texttt{if x < 5:}\\
    \hspace*{0.40in}\texttt{return x * 2}\\
    \hspace*{0.20in}\texttt{return x + 10}
  \end{block}

  \begin{activity}
    Without running Python, determine \texttt{transform(-2)}, \texttt{transform(4)}, and \texttt{transform(8)}. For each call, list the exact sequence of function calls, conditions, and returns.
  \end{activity}
\end{frame}

\begin{frame}{Takeaways}
  \begin{itemize}
    \item A conditional selects an execution path; \texttt{return} terminates the active function call.
    \item Several \texttt{return} statements represent alternative exits, not several results from one call.
    \item Chained conditionals select the first applicable return path; nested conditionals represent dependent decisions.
    \item Early returns and guard clauses can simplify control flow when a result is already known.
    \item Predicate calls must be evaluated to Boolean values before their surrounding conditions can be resolved.
    \item A complete trace records parameter bindings, evaluated conditions, selected branches, return values, and the caller's resumed state.
  \end{itemize}
\end{frame}

\end{document}
